\documentclass[a4paper, 12pt, oneside]{book}
\usepackage[utf8]{inputenc}
\usepackage[margin=3cm, bindingoffset=1cm]{geometry}
\linespread{1.5}
\usepackage{float}
\usepackage{csquotes}
\usepackage{subfig}
\usepackage{graphicx}
\usepackage{siunitx}
\usepackage{tikz}
\usepackage{xcolor}
\usepackage{titlesec}
\usepackage{fancyhdr}
\usepackage{parskip}
\usetikzlibrary{shapes.geometric, arrows.meta, positioning}
\usepackage[italian]{babel}

\pagestyle{fancy}
\setlength{\headheight}{14.5pt}
\renewcommand{\chaptermark}[1]{\markboth{\thechapter.\ \uppercase{#1}}{}}
\fancyhf{}
\fancyhead[L]{\small\nouppercase{\leftmark}}
\fancyhead[R]{}
\fancyfoot[C]{\thepage}
\renewcommand{\headrulewidth}{0.4pt}
\renewcommand{\footrulewidth}{1pt}

\titleformat{\chapter}[display]
{\normalfont\raggedright}
{\large\bfseries \MakeUppercase{\chaptertitlename\ \thechapter}}
{6pt}
{\Huge\bfseries}
[\vspace{10pt}\titlerule]

\titlespacing*{\chapter}{0pt}{10pt}{35pt}

\title{Tesi Triennale}
\author{Francesco Coppola}
\date{Ottobre 2026}

\newenvironment{dedication}
  {\clearpage           
   \thispagestyle{empty}
   \vspace*{\stretch{1}}
   \itshape            
   \raggedleft     
  }
  {\par 
   \vspace{\stretch{3}} 
   \clearpage           
  }
  
\begin{document}
%Frontespizio
\begin{titlepage}
    \begin{center}
        \LARGE{\uppercase{Università degli Studi di Salerno}}\\
        \vspace{5mm}
        %Dipartimento
    	\uppercase{\normalsize Dipartimento di Informatica}\\
    	%Corso di Laurea
    	\normalsize{ Corso di Laurea Triennale in Informatica}\\
    	\vspace{15mm}
    \end{center}
    \begin{figure}[H]
        \centering
        \includegraphics[width=0.35\textwidth]{logo_unisa}
    \end{figure}
    
    \begin{center}
    	%Titolo
		{\LARGE{\bf TESI DI LAUREA}}\\
		\vspace{3mm}
    	%Titolo
        {\LARGE{\bf Profili Clinici e Metabolici per l'Individuazione del Danno Renale Occulto: Modelli di Machine Learning e Stratificazione KDIGO}}\\
    	\vspace{3mm}
    \end{center}
    
    \vspace{15mm}
    \noindent
    \begin{minipage}[t]{0.47\textwidth}
        %Relatore
    	{\large{ Relatori:\\\bf Dott.ssa Fabiola De Marco\\Dott. Domenico Rossi}}
    	\vspace{12mm}\\
    \end{minipage}
    \hfill
    \begin{minipage}[t]{0.4\textwidth}\raggedleft
        %Candidato
    	{\large{Candidato: \\ \bf Francesco Coppola\\ Mat. 0512119837}}
    \end{minipage}
    
    \vspace{20mm}
    
    %Anno Accademico
    \centering{\large \uppercase{ Anno Accademico 2026/2027 }}

\end{titlepage}
%\begin{dedication}
%Dedicata a coloro che\\ ci provano ogni giorno.
%\end{dedication}
%Corpo della Tesi
\tableofcontents
\clearpage
\sloppy
\hyphenpenalty=10000
\exhyphenpenalty=10000
\chapter{Introduzione}
La malattia renale cronica è una condizione diffusa e spesso non riconosciuta: la maggior parte delle persone che ne sono  affette non sa di averla.~\cite{centers_for_disease_and_prevention_chronic_2026,wang_prevalence_2023}.

Riconoscere tale malattia non richiede esami complessi: bastano delle analisi del sangue e una delle urine~\cite{centers_for_disease_and_prevention_chronic_2026}. Eppure, nella pratica vengono prescritti soprattutto a chi ha già un motivo preciso per farli, come il diabete~\cite{stevens_kdigo_2024,centers_for_disease_and_prevention_chronic_2026,american_diabetes_association_professional_practice_committee_11_2024}. Questa tesi parte da qui, si chiede che cosa possa dire dei reni un modello di machine learning che dispone soltanto di dati raccolti in una normale visita, senza alcun esame renale. Il capitolo presenta il problema, gli approcci esistenti e i loro limiti. Seguono gli obiettivi, il metodo e i risultati principali.

\section{La malattia renale cronica}
Prima di parlare di modelli conviene chiarire come la malattia viene definita. Da questa definizione dipende tutto il resto del lavoro.

Secondo le linee guida KDIGO (\textit{Kidney Disease: Improving Global Outcomes}), la malattia renale cronica (CKD, Chronic Kidney Disease) è un' alterazione della struttura o della funzione dei reni che persiste per almeno tre mesi~\cite{stevens_kdigo_2024}. Nella pratica si riconosce con due esami. il primo è la creatinina nel sangue, da cui si stima la velocità di filtrazione glomerulare (\textbf{eGFR}, \textit{estimated Glomerular Filtration Rate}), cioè quanto plasma i reni riescono a filtrare in un minuto. Il secondo esame è il rapporto fra albumina e creatinina nelle urine (\textbf{ACR}, \textit{Albumin-to-Creatinine Ratio}), che misura quanta albumina il rene lascia passare. Sono considerati marcatori di malattia un eGFR $< \SI{60}{\milli\liter\per\minute\per1.73\meter\squared}$ oppure un ACR pari o superiore a $30\text{mg/g}$~\cite{stevens_kdigo_2024}. Incrociando i due valori, KDIGO divide il rischio in quattro livelli: basso, moderato, alto e molto alto ~\cite{stevens_kdigo_2024}.

La malattia è frequente e poco conosciuta da chi ne è affetto. Negli Stati Uniti la grande maggioranza degli adulti con CKD non sa di averla ~\cite{centers_for_disease_and_prevention_chronic_2026}. Lo stesso accade in Cina. Qui un'indagine nazionale ha trovato una consapevolezza molto bassa e ha indicato nell'albuminuria il marcatore più frequente, più della filtrazione ridotta ~\cite{wang_prevalence_2023}.

Per questo le linee guida indicano di cercare la malattia nelle persone a rischio con entrambi gli esami~\cite{stevens_kdigo_2024}. Per le persone con diabete è prevista la somministrazione annuale di ACR e eGFR, in tutto il diabete di tipo 2 e nel tipo 1 dopo cinque anni dalla diagnosi~\cite{american_diabetes_association_professional_practice_committee_11_2024}. La malattia renale cronica attribuita al diabete riguarda infatti fra il 20\% e il 40\% dei diabetici~\cite{american_diabetes_association_professional_practice_committee_11_2024}. Per chi non ha il diabete manca un'indicazione analoga: lo screening annuale dell'albuminuria è raccomandato solo nelle persone con diabete~\cite{bragg-gresham_using_2025}.

Testare tutti, d'altra parte, non è la soluzione più efficiente. Una revisione di 21 valutazioni economiche ha trovato lo screening più conveniente nelle popolazioni selezionate con punteggi di rischio più alti e non conveniente nelle popolazioni senza diabete né ipertensione~\cite{yeo_cost-effectiveness_2024}. Già nel 2003 Boulware e colleghi avevano trovato lo screening della proteinuria non conveniente, a meno di rivolgerlo ai gruppi a rischio più alto, come anziani e ipertesi, o di ripeterlo solo ogni dieci anni~\cite{boulware_screening_2003}.

Resta una domanda pratica, a chi prescrivere gli esami renali?

\section{Approcci presenti in letteratura}
Per decidere chi sottoporre a un esame serve una stima del rischio della singola persona. in medicina questa stima viene da modelli di predizione clinica, che stimano la probabilità che una persona presenti una condizione (modelli diagnostici) o che la sviluppi in futuro (modelli prognostici)~\cite{collins_tripodai_2024}.

Con la digitalizzazione dei dati sanitari, le cartelle cliniche elettroniche permettono di addestrare un modello su decine di milioni di cartelle di pazienti, con centinaia di miliardi di dati~\cite{rajkomar_machine_2019}. A differenza dei sistemi tradizionali, programmati con regole esplicite, nel machine learning il modello impara dagli esempi~\cite{rajkomar_machine_2019}. La diffusione di questi metodi in medicina è stata resa possibile da grandi insiemi di dati etichettati~\cite{topol_high-performance_2019}, e con l'accesso crescente ai dati e a software pronti all'uso sviluppare un modello di predizione è diventato più rapido e semplice~\cite{collins_tripodai_2024}. Anche per la malattia renale esiste un ampia gamma di modelli di machine learning~\cite{sanmarchi_predict_2023}.

Questi modelli possono essere raggruppati in tre filoni. il primo comprende i punteggi clinici per individuare le persone a rischio nella popolazione generale o nelle cure primarie. SCORED stima la probabilità di un eGFR ridotto a partire da nove informazioni semplici, fra cui età, anemia, diabete e proteinuria~\cite{bang_screening_2007}. QKidney stima il rischio di sviluppare la malattia nei cinque anni successivi, usando i dati delle cure primarie di Inghilterra e Galles~\cite{hippisley-cox_predicting_2010}.

Il secondo filone riguarda l'albuminuria, il marcatore più difficile da stimare senza un esame delle urine. Bragg-Gresham e colleghi hanno costruito un modello per stabilire quali adulti statunitensi senza diabete sottoporre allo screening dell'albuminuria~\cite{bragg-gresham_using_2025}. MERWACS usa il machine learning per individuare, fra gli adulti con almeno 50 anni, chi sottoporre per primo agli esami di laboratorio per la malattia renale~\cite{yoo_merwacs_2026}.

Il terzo filone lavora a valle della diagnosi. La \textit{Kidney Failure Risk Equation}, validata su oltre 700.000 pazienti di più di 30 paesi, stima nei pazienti con malattia già accertata il rischio di arrivare alla dialisi o al trapianto, e parte proprio da eGFR e ACR~\cite{tangri_multinational_2016}.

Il quarto filone è l'apprendimento non supervisionato, che cerca gruppi o anomalie nei dati senza usare un'etichetta~\cite{hastie_elements_2009,beam_big_2018}. Nella malattia renale è stato usato soprattutto per descrivere sottogruppi di pazienti già diagnosticati~\cite{zheng_subtyping_2021}.

Questo strumento si colloca fra il primo e il secondo filone: stima i marcatori presenti oggi, prima che gli esami renali vengano eseguiti.

\section{Limiti degli approcci esistenti e obiettivi della tesi}
La letteratura appena descritta lascia aperti alcuni problemi, che conviene raccogliere in un unico punto perchè motivano le scelte di questo lavoro.

Il primo limite riguarda gli input, Alcuni modelli includono fra i predittori un esame che misura direttamente il danno renale: SCORED, per esempio, include la proteinuria~\cite{bang_screening_2007}. Nel caso estremo il modello ricopia la definizione della malattia. Uno dei dataset pubblici dell'UCI (\textit{University of California, Irvine}) sulla malattia renale cronica contiene creatinina e urea del sangue, albumina e peso specifico delle urine~\cite{l_rubini_chronic_2015}, e due lavori che lo usano riportano accuratezze altissime fra il 98,9\% e il 99,8\%~\cite{qin_machine_2020,chittora_prediction_2021}. Prestazioni simili sono ottime per misurare la regola diagnostica non la capacità di anticiparla.

Il secondo limite riguarda il \textit{target}. Prestazioni più alte si ottengono quando si stima l'eGFR ridotto~\cite{bang_screening_2007,sabanayagam_deep_2020}, che per costruzione dipende dall'età, inclusa nell'equazione che lo calcola~\cite{inker_new_2021}. Quando il \textit{target} è l'albuminuria, il marcatore più frequente nella popolazione~\cite{wang_prevalence_2023}, i modelli senza esami delle urine ottengono prestazioni più modeste~\cite{bragg-gresham_using_2025,muntner_development_2011}. Quando i due criteri sono riuniti in un'unica etichetta, non si vede più quale dei due il modello riconosca~\cite{ramspek_prediction_2021, song_using_2020}.

Il terzo limite riguarda la validazione. Una revisione sistematica conta i modelli di rischio per la malattia renale, di cui solo una parte era stata validata su dati esterni~\cite{echouffo-tcheugui_risk_2012}. Haris e colleghi segnalano un alto rischio di distorsione in molti di essi~\cite{haris_prediction_2024}. I modelli di machine learning, inoltre, vengono raramente verificati su popolazioni diverse da quella di sviluppo~\cite{sanmarchi_predict_2023}. Non è nemmeno scontato che la complessità del modello utilizzato serva, in un confronto con la regressione logisitica una regressione sistematica non ha trovato un vantaggio nel machine learning~\cite{christodoulou_systematic_2019}.

C'è poi la soglia. KDIGO separa sani e malati con un valore fisso~\cite{stevens_kdigo_2024}. Il rischio legato all'albuminuria, però, non onosce quel confine e cresce in modo continuo anche al di sotto~\cite{matsushita_association_2010}. Un modello addestrato sull'etichetta eredita la soglia. Di ciò che accade prima non può dire nulla.

L'ultimo punto riguarda l'utilità clinica. Nessuno dei modelli esaminati da Haris e colleghi presenta analisi dell'utilità clinica o studi di impatto che dicano cosa cambierebbe usandoli nella pratica~\cite{haris_prediction_2024}. Van Calster e colleghi raccomandano invece di riportare sempre una misura di questo genere~\cite{van_calster_evaluation_2025}.

\section{Obiettivi della tesi}
Da questi limiti discende l'obiettivo del lavoro. La tesi vuole stabilire che cosa un modello riesce a riconoscere della malattia renale senza vedere alcun esame renale. Vuole anche approfondire se ciò che il modello impara resta valido su un'altra popolazione e ha dell'utilità clinica. Gli obiettivi specifici sono cinque.
\begin{enumerate}
	\item \textbf{Misurare la capcità predittiva senza \textit{leakage}.} Stimare i marcatori KDIGO a partire dai soli dati di routine e valutare i due criteri anche separatamente.
	\item \textbf{Misurare se il modello riconosce la stratificazione in livelli di KDIGO} Stimare se i casi individuati dal modello rientrano nelle fasce di gravità identificate da KDIGO.
	\item \textbf{Confrontare un modello senza etichette con KDIGO.} Applicare l'apprendimento non supervisionato, senza formula e soglie, e osservare le sue performance e come si comporta rispetto all'apprendimento con etichette e alla definizione clinica.
	\item \textbf{Confrontare il modello con la pratica.} Verificare se, a parità di esami prescritti, il modello individua più casi del criterio basato su diabete e  ipertensione.
	\item \textbf{Verificare il trasporto del modello su un'altra popolazione}. Applicare i modelli, senza riaddestrarli, a una popolazione esterna a quella di sviluppo.
\end{enumerate}

\noindent Il lavoro non propone un nuovo algoritmo e non costruisce uno strumento diagnostico. Con una sola misurazione per persona non si può verificare la persistenza richiesta dalla definizione~\cite{stevens_kdigo_2024}: per questo si parla di marcatori e non di diagnosi.

\section{Approccio metodologico}
Per raggiungiere gli obiettivi prefissati è stata costruita una procedura che parte dai dati grezzi e arriva alla valutazione finale, schematizzata nella figura~\ref{fig:schema_metodologico}.

I dati provengono da un dataset pubblico~\cite{li_bimodal_2026} raccolto nel reparto di Diabetologia ed Endocrinologia di un ospedale di Shanghai sulla minoranza Jinuo~\cite{jiang_jinuo__2016}. Per ogni persona l'etichetta è costruita dagli esami renali secondo KDIGO~\cite{stevens_kdigo_2024}, con la velocità di filtrazione glomerulare ricalcolata con un'equazione validata~\cite{inker_new_2021}. Dopo questo passaggio gli esami renali poi vengono esclusi dall'input. Prima di qualsiasi stima che impari dai dati, una parte dei soggetti viene messa da parte come insieme di test: resta chiusa per tutto lo sviluppo e viene letta una sola volta, alla fine, per confermare le conclusioni.

Sui dati restanti il lavoro si divide in tre fasi. Si confrontano anzitutto modelli supervisionati di complessità crescente, dalla regressione logistica agli alberi decisionali~\cite{breiman_random_2001,chen_xgboost_2016}, in validazione incrociata annidata. Lo stesso confronto è ripetuto sui due criteri della regola KDIGO presi uno alla volta. Poi le persone vengono ordinate per rischio stimati e si contano i casi trovati a parità di esami prescritti. In seguito si addestrano i modelli non supervisionati, con e senza le variabili renali, e il loro punteggio viene messo a confronto con la definizione KDIGO.

Infine, una volta stimati i modelli, vengono applicati così come sono a un'indagine nazionale statunitense~\cite{nhanes_cdc}. Nessun parametro viene ristimato. Poichè di questa popolazione è nota anche la mortalità negli anni successivi sono state effettuate ulteriori analisi per stimare se il punteggio non supervisionato ha un significato clinico.

\begin{center}
\begin{tikzpicture}[
	node distance=0.6cm,
	block/.style={rectangle, draw=blue!50!black, text width=0.75\textwidth, text centered, rounded corners=4pt, minimum height=0.9cm, fill=blue!5, thick, font=\sffamily},
	line/.style={draw, -Stealth, thick, gray!60!black}
	]
	
	\node [block] (n1) {Dataset pubblico (coorte ospedaliera)};
	\node [block, below=of n1] (n2) {Etichetta: marcatori KDIGO};
	\node [block, below=of n2] (n3) {Separazione dell'insieme di test};
	\node [block, below=of n3] (n4) {Esclusione degli esami renali};
	\node [block, below=of n4] (n5) {Modelli di complessità crescente};
	\node [block, below=of n5] (n6) {Bilanciamento delle classi e sottogruppo diabetico};
	\node [block, below=of n6] (n7) {Confronto con i livelli KDIGO};
	\node [block, below=of n7] (n8) {Apprendimento non supervisionato};
	\node [block, below=of n8] (n9) {Conferma sull'insieme di test};
	\node [block, below=of n9] (n10) {Validazione su coorte esterna};
	
	\path [line] (n1) -- (n2);
	\path [line] (n2) -- (n3);
	\path [line] (n3) -- (n4);
	\path [line] (n4) -- (n5);
	\path [line] (n5) -- (n6);
	\path [line] (n6) -- (n7);
	\path [line] (n7) -- (n8);
	\path [line] (n8) -- (n9);
	\path [line] (n9) -- (n10);
	
\end{tikzpicture}
\captionof{figure}{Schema generale della procedura. Ogni blocco corrisponde a un momento del lavoro descritto nel testo; procedure e parametri sono riportati nel Capitolo 3.}
\label{fig:schema_metodologico}
\end{center}

\section{Risultati ottenuti}
I risultati sono riportati per esteso nel capitolo 4. Qui se ne dà una sintesi.

Senza alcun esame renale, il modello con le prestazioni più alte, il Random Forest, distingue le persone con marcatori di malattia renale con un'area sotto la curva ROC (AUROC) di 0,703 in validazione incrociata. Sull'insieme di test, mai usato durante lo sviluppo, l'AUROC è di 0,743. Una trentina di modelli diversi non va oltre. La ragione sta nei due criteri. Dal sangue la filtrazione ridotta si riconosce bene, con un'AUROC di 0,846, mentre l'albuminuria, che riguarda nove casi positivi su dieci, si ferma a 0,690. Il risultato è in linea con i modelli pubblicati per un \textit{target} simile~\cite{yoo_merwacs_2026,muntner_development_2011}. Undici strategie di bilanciamento dati aggiuntive, compresi modelli più recenti, non migliorano il risultato in modo apprezzabile, e nemmeno l'addestramento del modello sul valore continuo dell'albuminuria invece che sulla sola soglia diagnostica.

Sul piano metodologico, la tesi mette alla prova una pratica sempre più usata, la correzione dello sbilanciamento delle classi~\cite{van_den_goorbergh_harm_2022}. A parità di prestazioni, nessuna tecnica fa riconoscere più casi gravi. 

Il rischio stimato rispetta la gravità clinica. In tutti i modelli la probabilità media cresce a ogni livello KDIGO, e fra il 67\% il 70\% delle coppie di persone di livello diverso è ordinato come nella stratificazione clinica. Con una soglia fissata per trovare il 90\% dei casi, i modelli riconoscono fra 46 e 47 dei 50 casi di rischio alto e fra 18 e 20 dei 21 casi di rischio molto alto.

A parità di persone inviate ad esame, il modello trova 120 casi. Il criterio basato su diabete e ipertensione ne trova 98, cioè 22 in meno (intervallo di confidenza al 95\% della differenza da 6 a 38) Lo stesso modello è stato poi applicato, senza riaddestramento, a 34.494 adulti statunitensi. Ne individua 4.710 contro 4.300, con una differenza di 410 casi (da 351 a 469).

Quanto al modello non supervisionato, ritrova la regola KDIGO solo se le variabili renali sono fra i dati, e in quel caso l'AUROC è 0,99. Senza esami renali scende a 0,65. È quanto ottiene la sola età. il suo punteggio, però. cresce con l'albuminuria già al di sotto della soglia e non mostra alcun salto in corrispondenza di essa. Fra 28.433 persone statunitentsi sane secondo KDIGO, una deviazione standard in più nel punteggio corrisponde a un rischio di morte cardiovascolare 1,38 volte maggiore (da 1,20 a1,58). Con le morti per incidente, usate come controllo, l'associazione non c'è.

Per la filtrazione ridotta i modelli hanno indicato una variabile rilevante, l'acido urico. Aggiunto a un punteggio semplice di cinque variabili, come quello utilizzato da SCORED~\cite{bang_screening_2007}, porta l'AUROC sulla popolazione statunitense da 0,886 a 0,914. Il modello più complesso, trasportato allo stesso modo, si ferma a 0,832.

Tutte le conclusioni sono state sottoposte all'insieme di test, aperto una sola volta alla fine.

\section{Struttura della tesi}
Il resto dell'elaborato è organizzato in quattro capitoli.

\textbf{Capitolo 2, Stato dell'arte.} La malattia renale cronica e i criteri KDIGO, le strategie di screening, i modelli di rischio con e senza esami renali, il machine learning supervisionato e non supervisionato, il problema delle classi sbilanciate, il \textit{leakage} e le misure per valutare un modello clinico anche su coorte esterna.

\textbf{Capitolo 3, Metodologia.} Il dataset e la costruzione dell'etichetta, la selezione delle variabili e il trattamento dei valori mancanti, i modelli e il protocollo di validazione, le tecniche di bilanciamento, le analisi esplorative, i modelli non supervisionati, la conferma sull'insieme di test, la validazione esterna e il prototipo interattivo.

\textbf{Capitolo 4, Risultati ottenuti.} I risultati sui marcatori KDIGO e sui due criteri separati, il confronto con la pratica, il confronto fra modello non supervisionato e KDIGO, la verifica sulla seconda popolazione e il confronto con lo stato dell'arte.

\textbf{Capitolo 5, Conclusioni e sviluppi futuri.} I contributi, i limiti del lavoro e le possibili estensioni.

\chapter{Stato dell'arte}

\chapter{Metodologia}

\chapter{Risultati ottenuti}

\chapter{Conclusioni e sviluppi futuri}

\chapter*{Ringraziamenti}
\addcontentsline{toc}{chapter}{Ringraziamenti}
\clearpage
\cleardoublepage
\addcontentsline{toc}{chapter}{Bibliografia}
\bibliographystyle{unsrt} 
\bibliography{bib}       
\clearpage
\listoffigures
\addcontentsline{toc}{chapter}{\listfigurename}
\clearpage
\listoftables
\addcontentsline{toc}{chapter}{\listtablename}
\clearpage
\end{document}
